\section{Flow and matching}

\entry{Dinic (max flow, iterative)}{$O(n^2 m)$, unit caps $O(m\sqrt n)$ \quad mem $n + m$}{%
Directed, 0-indexed. No recursion, so path length never matters; multi-edges,
self-loops and 0-capacity edges are fine. \code{add(u, v, c)} returns an \textbf{even}
edge id, \code{id \^{} 1} is its reverse; an \textbf{undirected} edge is
\code{add(u, v, c, c)}.\\[0.3ex]
\begin{ctab}{0.22}{0.73}
\code{flow(id)} & flow on edge \code{id}; negative = it runs \code{v -> u}.\\
\code{inS(v)} & after \code{maxflow}: \code{v} on the \code{s} side of the min cut
closest to \code{s}.\\
\code{cut()} & edge ids crossing $S \to T$, all saturated; their \code{c} sums to the
flow. An odd id = undirected edge crossing reversed, \code{id \& \~{}1} is what
\code{add} returned.\\
\code{reset()} & all flow back to 0, edges kept.\\
\end{ctab}\\[0.3ex]
\textbf{Config.} \code{F} = capacity type (\code{int} saves memory, \code{double} works
with \code{EPS = 1e-9}). \code{INF} is max/4, so up to 3 can add up; a result
\code{>= INF} means no finite cut.\\[0.3ex]
\code{maxflow} is \textbf{incremental}: add edges or raise \code{d.e[id].c}, call it
again, and it returns only the \emph{extra} flow. A different \code{(s, t)} also
continues from the old flow --- \code{reset()} first.\\[0.3ex]
\textbf{Recipes.} Bipartite matching: use \code{Matching} below, about 5x faster than
the same graph here. Min vertex cover in a bipartite graph = max matching (K\"onig), and
max independent set = $n$ minus that.
Vertex capacities: split \code{v} into \code{v\_in -> v\_out} with that capacity.
Minimum path cover of a DAG = $n$ minus max matching on the split graph.
Project selection / max closure: source$\to$profit nodes, cost nodes$\to$sink,
prerequisites as \code{INF} edges, answer = total profit $-$ max flow.\\[0.3ex]
Measured at -O2: random $10^5$ vertices $10^6$ edges 0.4\,s; unit-cap bipartite
$2\cdot10^5$ / $1.2\cdot10^6$ 0.8\,s; a $300\times300$ grid with random two-way caps
4.5\,s --- a bad case for every Dinic.}

%LABEL Dinic
\begin{lstlisting}
struct Dinic {
    using F = ll;
    static constexpr F EPS = 0;            // double: 1e-9
    static constexpr F INF = numeric_limits<F>::max() / 4;
    struct E { int to; F c, f; };          // cap, flow
    int n; vector<E> e; vector<vector<int>> g;
    vector<int> lv, it;                    // level, edge ptr

    Dinic(int n = 0) : n(n), g(n), lv(n), it(n) {}
    int add(int u, int v, F c, F rc = 0) {     // id of u->v
        g[u].push_back(e.size()); e.push_back({v, c, 0});
        g[v].push_back(e.size()); e.push_back({u, rc, 0});
        return (int)e.size() - 2;
    }
    F res(int id) { return e[id].c - e[id].f; }
    F flow(int id) { return e[id].f; }
    bool inS(int v) { return lv[v] >= 0; }
    void reset() { for (auto& x : e) x.f = 0; }

    bool bfs(int s, int t) {
        fill(lv.begin(), lv.end(), -1); lv[s] = 0;
        vector<int> q{s};
        for (int i = 0; i < (int)q.size() && lv[t] < 0; i++)
            for (int id : g[q[i]]) {
                int v = e[id].to;
                if (lv[v] < 0 && res(id) > EPS)
                    lv[v] = lv[q[i]] + 1, q.push_back(v);
            }
        return lv[t] >= 0;
    }
    F maxflow(int s, int t) {
        F tot = 0; vector<int> st;         // path, edge ids
        while (bfs(s, t)) {
            fill(it.begin(), it.end(), 0);
            int u = s;
            while (true) {
                if (u == t) {                  // augment
                    F b = INF; int k = 0;  // k = 1st sat.
                    for (int id : st) b = min(b, res(id));
                    for (int i = st.size(); i--;) {
                        e[st[i]].f += b; e[st[i] ^ 1].f -= b;
                        if (res(st[i]) <= EPS) k = i;
                    }
                    tot += b; st.resize(k);
                    u = k ? e[st.back()].to : s;
                } else if (it[u] < (int)g[u].size()) {
                    int id = g[u][it[u]], v = e[id].to;
                    if (lv[v] == lv[u] + 1 && res(id) > EPS)
                        st.push_back(id), u = v;
                    else it[u]++;
                } else {                       // dead end
                    if (u == s) break;
                    lv[u] = -1;
                    u = e[st.back() ^ 1].to;
                    st.pop_back(); it[u]++;
                }
            }
        }
        return tot;
    }
    vector<int> cut() {                    // after maxflow
        vector<int> r;
        for (int id = 0; id < (int)e.size(); id++) {
            int u = e[id ^ 1].to, v = e[id].to;
            if (e[id].c > 0 && inS(u) && !inS(v))
                r.push_back(id);
        }
        return r;
    }
};
\end{lstlisting}

\entry{Min cost max flow}{$O(\mathrm{aug} \cdot m \log n)$ \quad mem $n + m$}{%
Dijkstra with potentials. \textbf{Negative costs are fine} (one SPFA pass first);
negative \emph{cycles} are not. No recursion. \code{add(u, v, cap, cost)} returns the
edge id, cost is per unit; \code{flow(id)} reads it back.\\[0.3ex]
\begin{ctab}{0.36}{0.59}
\code{run(s, t)} & \code{\{flow, cost\}}: max flow, min cost among those.\\
\code{run(s, t, k)} & at most \code{k} units, min cost.\\
\code{run(s, t, m.INFF, 1)} & \emph{any} amount, min cost: stops once one more unit
costs \code{>= 0}.\\
\end{ctab}\\[0.3ex]
\textbf{Config.} \code{F} flow type, \code{C} cost type (\code{double} works --- compare
with a tolerance). \code{INFF} / \code{INFC} are max/4.\\[0.3ex]
\textbf{Max} profit: add every edge with \code{-cost}, negate the answer. Undirected edge
with cost: two separate \code{add} calls. Paths go cheapest first, so after $k$ units
the cost is the minimum for exactly $k$. \code{run} is incremental like \code{Dinic}.\\[0.3ex]
One Dijkstra per augmenting path ($\le$ total flow), fine up to $\sim10^4$--$10^5$ of
them. Measured: $10^4$ vertices $10^5$ edges 0.4\,s; a $300\times300$ assignment
0.42\,s, while \code{hungarian} does $1000\times1000$ in 0.08\,s --- use that for a
dense square cost matrix.}

%LABEL MCMF
\begin{lstlisting}
struct MCMF {
    using F = ll; using C = ll;
    static constexpr F INFF = numeric_limits<F>::max() / 4;
    static constexpr C INFC = numeric_limits<C>::max() / 4;
    struct E { int to; F c, f; C w; };     // cap, flow, cost
    int n; vector<E> e; vector<vector<int>> g;
    vector<C> pi, dist; vector<int> pe;    // pe = parent edge

    MCMF(int n = 0) : n(n), g(n), pi(n, 0), dist(n), pe(n) {}
    int add(int u, int v, F c, C w) {
        g[u].push_back(e.size()); e.push_back({v, c, 0, w});
        g[v].push_back(e.size()); e.push_back({u, 0, 0, -w});
        return (int)e.size() - 2;
    }
    F flow(int id) { return e[id].f; }

    void spfa(int s) {                     // exact potentials
        fill(pi.begin(), pi.end(), INFC); pi[s] = 0;
        vector<char> inq(n, 0); deque<int> q{s};
        while (!q.empty()) {
            int u = q.front(); q.pop_front(); inq[u] = 0;
            for (int id : g[u]) {
                auto& x = e[id];
                if (x.c - x.f > 0 && pi[u] + x.w < pi[x.to]) {
                    pi[x.to] = pi[u] + x.w;
                    if (!inq[x.to])
                        inq[x.to] = 1, q.push_back(x.to);
                }
            }
        }
        for (auto& p : pi) if (p == INFC) p = 0;  // unreached
    }
    bool dijkstra(int s, int t) {
        fill(dist.begin(), dist.end(), INFC); dist[s] = 0;
        priority_queue<pair<C,int>, vector<pair<C,int>>,
                       greater<pair<C,int>>> pq;
        pq.push({0, s});
        while (!pq.empty()) {
            auto [d, u] = pq.top(); pq.pop();
            if (d > dist[u]) continue;
            for (int id : g[u]) {
                auto& x = e[id];
                C nd = d + x.w + pi[u] - pi[x.to];   // >= 0
                if (x.c - x.f > 0 && nd < dist[x.to]) {
                    dist[x.to] = nd; pe[x.to] = id;
                    pq.push({nd, x.to});
                }
            }
        }
        for (int v = 0; v < n; v++)
            if (dist[v] < INFC) pi[v] += dist[v];
        return dist[t] < INFC;
    }
    pair<F,C> run(int s, int t, F lim = INFF, bool neg = 0) {
        F fl = 0; C cost = 0;
        for (auto& x : e)
            if (x.c - x.f > 0 && x.w < 0) { spfa(s); break; }
        while (fl < lim && dijkstra(s, t)) {
            C pc = pi[t] - pi[s];          // cost of one unit
            if (neg && pc >= 0) break;
            F b = lim - fl;
            for (int v = t; v != s; v = e[pe[v] ^ 1].to)
                b = min(b, e[pe[v]].c - e[pe[v]].f);
            for (int v = t; v != s; v = e[pe[v] ^ 1].to)
                e[pe[v]].f += b, e[pe[v] ^ 1].f -= b;
            fl += b; cost += b * pc;
        }
        return {fl, cost};
    }
};
\end{lstlisting}

\entry{Bipartite matching (Hopcroft-Karp)}{$O(m\sqrt n)$ \quad mem $n + m$}{%
Unweighted maximum matching, iterative, multi-edges fine. Left \code{[0, nl)} and right
\code{[0, nr)} are numbered \textbf{independently}; \code{add(l, r)}, then
\code{run()} returns the size $k$. \code{ml[l]} / \code{mr[r]} give the partner or
\code{-1}.\\[0.3ex]
\textbf{K\"onig.} After \code{run}, \code{auto [L, R] = m.cover()} is a minimum vertex
cover, \code{|L| + |R| == k}. Its complement --- left \emph{not} in \code{L} plus right
\emph{not} in \code{R} --- is a maximum independent set, size \code{nl + nr - k}.\\[0.3ex]
\code{run} is incremental: add edges and call it again, old pairs are kept. You may also
pre-fill \code{ml}/\code{mr} with a valid greedy matching.\\[0.3ex]
Weighted: \code{hungarian} or \code{MCMF}. General (non-bipartite) graphs need blossom,
not this. Measured: $10^5 + 10^5$ vertices, $10^6$ edges 0.15\,s, 15\,MB.}

%LABEL Matching
\begin{lstlisting}
struct Matching {
    int nl, nr; vector<vector<int>> g;
    vector<int> ml, mr, d, it;             // d = bfs layer

    Matching(int nl, int nr)
        : nl(nl), nr(nr), g(nl), ml(nl, -1), mr(nr, -1) {}
    void add(int l, int r) { g[l].push_back(r); }

    bool bfs() {                           // layers
        d.assign(nl, -1); vector<int> q; bool found = false;
        for (int l = 0; l < nl; l++)
            if (ml[l] < 0) d[l] = 0, q.push_back(l);
        for (int i = 0; i < (int)q.size(); i++)
            for (int r : g[q[i]]) {
                int w = mr[r];
                if (w < 0) found = true;
                else if (d[w] < 0)
                    d[w] = d[q[i]] + 1, q.push_back(w);
            }
        return found;
    }
    bool aug(int s) {                      // s = free left
        vector<int> st{s};
        while (!st.empty()) {
            int u = st.back();
            if (it[u] == (int)g[u].size()) {
                d[u] = -1; st.pop_back(); continue;
            }
            int r = g[u][it[u]], w = mr[r];
            if (w < 0) {                   // free: flip path
                for (int x : st)
                    ml[x] = g[x][it[x]], mr[ml[x]] = x;
                return true;
            }
            if (d[w] == d[u] + 1) st.push_back(w);
            else it[u]++;
        }
        return false;
    }
    int run() {
        while (bfs()) {
            it.assign(nl, 0);
            for (int l = 0; l < nl; l++)
                if (ml[l] < 0) aug(l);
        }
        int k = 0;
        for (int l = 0; l < nl; l++) k += ml[l] >= 0;
        return k;
    }
    pair<vector<int>, vector<int>> cover() { // after run
        vector<char> vl(nl, 0), vr(nr, 0);
        vector<int> q, L, R;
        for (int l = 0; l < nl; l++)
            if (ml[l] < 0) vl[l] = 1, q.push_back(l);
        for (int i = 0; i < (int)q.size(); i++)
            for (int r : g[q[i]]) if (!vr[r]) {
                vr[r] = 1; int w = mr[r];
                if (w >= 0 && !vl[w])
                    vl[w] = 1, q.push_back(w);
            }
        for (int l = 0; l < nl; l++)
            if (!vl[l]) L.push_back(l);
        for (int r = 0; r < nr; r++)
            if (vr[r]) R.push_back(r);
        return {L, R};
    }
};
\end{lstlisting}

\entry{Hungarian (assignment)}{$O(n^2 m)$ \quad mem $nm$}{%
Minimum-cost \textbf{perfect} matching on a cost matrix, rows $\le$ columns: every row
gets its own column, cheapest in total. Returns \code{\{cost, asg\}} with
\code{asg[row] = column}. Input is 0-indexed; the internals are 1-indexed, which is why
the matrix is read as \code{a[i0-1][j-1]}.\\[0.3ex]
\textbf{Maximise} instead: negate every cost, then negate the answer. Forbidden pairs:
a large finite cost, not \code{INF} (\code{INF} arithmetic overflows here) --- something
like $10^{12}$ with real costs below $10^6$.\\[0.3ex]
Not the same thing as \code{Matching}: that one maximises the \emph{number} of pairs,
this one fixes the number and minimises the \emph{cost}. Costs may be negative.}

%LABEL hungarian
\begin{lstlisting}
// a = cost matrix, a.size() <= a[0].size()
pair<ll, vector<int>> hungarian(const matrix& a) {
    int n = (int)a.size(), m = (int)a[0].size();
    // 1-indexed inside; column 0 is virtual, p[j]==0 => free
    vi u(n + 1, 0), v(m + 1, 0), minv(m + 1);
    vector<int> p(m + 1, 0), way(m + 1, 0);
    for (int i = 1; i <= n; i++) {
        p[0] = i; int j0 = 0;
        fill(all(minv), INF);
        vector<char> used(m + 1, 0);
        do {
            used[j0] = 1;
            int i0 = p[j0], j1 = 0; ll delta = INF;
            for (int j = 1; j <= m; j++) if (!used[j]) {
                ll cur = a[i0-1][j-1] - u[i0] - v[j];
                if (cur < minv[j]) {
                    minv[j] = cur; way[j] = j0;
                }
                if (minv[j] < delta) {
                    delta = minv[j]; j1 = j;
                }
            }
            for (int j = 0; j <= m; j++) {
                if (used[j]) {
                    u[p[j]] += delta; v[j] -= delta;
                } else minv[j] -= delta;
            }
            j0 = j1;
        } while (p[j0] != 0);
        do {                          // flip the path back
            int j1 = way[j0]; p[j0] = p[j1]; j0 = j1;
        } while (j0);
    }
    vector<int> asg(n);
    for (int j = 1; j <= m; j++) if (p[j]) asg[p[j]-1] = j-1;
    ll cost = 0;
    for (int i = 0; i < n; i++) cost += a[i][asg[i]];
    return {cost, asg};
}
\end{lstlisting}

\entry{Eulerian path}{$O(n + m)$ \quad mem $n + m$}{%
Hierholzer, iterative, handles both directions with the \code{dir} flag. Returns the
\textbf{vertex} sequence, length \code{m+1}, using every edge exactly once --- or an
empty vector when none exists (wrong degrees, or the edges are disconnected). A closed
tour comes back with \code{path.front() == path.back()}.\\[0.3ex]
Existence: \emph{directed} needs \code{out-in} equal to 0 everywhere except possibly one
$+1$ (the start) and one $-1$ (the end); \emph{undirected} needs zero or exactly two odd
vertices. Both also need all edges in one connected component --- which the final
\code{size() != m+1} test catches, so isolated vertices do not matter.\\[0.3ex]
For the edge sequence instead, record \code{id} alongside each pop.}

%LABEL eulerPath
\begin{lstlisting}
vector<int> eulerPath(int n, const vector<pii>& es,
                      bool dir) {
    int m = (int)es.size();
    if (!m) return {};
    vector<vector<pii>> g(n);          // (to, edge id)
    vi outd(n, 0), ind(n, 0);
    for (int i = 0; i < m; i++) {
        int a = (int)es[i].first, b = (int)es[i].second;
        g[a].push_back({b, i}); outd[a]++; ind[b]++;
        if (!dir) {
            g[b].push_back({a, i}); outd[b]++; ind[a]++;
        }
    }
    int s = -1;
    if (dir) {
        int plus = 0, minus = 0;
        for (int v = 0; v < n; v++) {
            ll df = outd[v] - ind[v];
            if (df == 1) { plus++; s = v; }
            else if (df == -1) minus++;
            else if (df) return {};
        }
        if (plus > 1 || minus > 1) return {};
    } else {
        int odd = 0;
        for (int v = 0; v < n; v++)
            if (outd[v] & 1) { odd++; s = v; }
        if (odd != 0 && odd != 2) return {};
    }
    if (s < 0)
        for (int v = 0; v < n; v++)
            if (outd[v]) { s = v; break; }
    vector<int> it(n, 0), st{s}, path;
    vector<char> used(m, 0);
    while (!st.empty()) {
        int u = st.back();
        while (it[u] < (int)g[u].size()
               && used[g[u][it[u]].second]) it[u]++;
        if (it[u] == (int)g[u].size()) {
            path.push_back(u); st.pop_back();
        } else {
            auto [v, id] = g[u][it[u]++];
            used[id] = 1; st.push_back((int)v);
        }
    }
    if ((int)path.size() != m + 1) return {};   // disjoint
    reverse(all(path));
    return path;
}
\end{lstlisting}
